\section{Discussion and Limitations}
The benchmark asks two questions of the same weights. The training objective measures how well a method optimized the observations it received. Error on new input values or sequence indices measures whether those weights continued the rule. On an unseen input no optimizer can infer a unique continuation from central labels alone. For example, for any real number $a$, the functions $0$ and $aR(x-3)$ agree on the entire central interval $[-3,3]$. At $x=10$ their predictions differ by $7a$. Knowledge of the task family, the architecture and Astra's pretraining can all affect which continuation a solver chooses.

The fixed-network comparison covers ten deterministic numerical tasks, with up to 50 trainable coefficients. The architecture-selection experiment additionally changes the allowed graph and cost. Continuous tasks have three sampled training sets. Integer tasks have one training set and three initializations. Multiple sizes in the same WSP response share context. We give counts over these fixed cases and show each size. Population claims about optimizers in vision, language or general large-scale training require other workloads, models and budgets. The iterative method uses its development set for both feedback and selection. Three attempts cap the cost, and a rejected proposal can still perform better on unseen inputs. The original numerical baselines see fewer labels; we therefore also report an intervention giving the two gradient procedures the entire feedback pool. Neither design matches computation or pretraining. The examples-only condition still discloses the architecture and intended range, and there is no common unit equating an Astra inference to a gradient or Jacobian evaluation. Each hosted-model response is archived because the provider model name alone cannot guarantee its reproduction.

\paragraph{Which sizes are known to suffice?}
For a one-state running sum, setting all three weights to one returns every sum in the same addition order as the reference. Autonomous Fibonacci and triangular generators require minimum real-arithmetic state widths two and three, respectively, for their infinite sequences (Appendix~\ref{app:additional} and the realization argument in Appendix~\ref{app:proof}). For arithmetic, write $t=x_1+x_2$ or $x_1-x_2$. Even a single hidden unit can return $R(t+20)-20=t$ throughout the bounded real domain. The large intermediate offset can change its binary64 rounding. Two units give $R(t)-R(-t)=t$ without that offset. For the shallow scalar architecture the best possible uniform real error on square over $[-10,10]$ is $50/(w+1)^2$ at width $w$ (Appendix~\ref{sec:tasks}). A different deep graph returned by Astra provably rounds every square correctly on the outer shell; an engineered extension covers the full binary64 interval (Appendices~\ref{app:literal-square-certificate} and~\ref{app:fullsquare}). Those architectures and proofs give constructive bounds. They do not establish a minimal deep network for the smooth functions.

The supplementary edit theorem fixes a single output coordinate and counts how many initialized coefficients must change to implement a recurrence (Appendix~\ref{sec:front}). That question has a different unit of cost from the primary comparison, where each solver can change every weight. The supplement also studies a register machine, an abstract computer that updates integer registers according to conditional instructions. Its proof concerns unbounded exact integers; the reported binary64 runs test bounded executions (Appendix~\ref{app:machine}). Neither construction is counted as a win in the optimizer comparison.

\paragraph{Interpreting aggregate error and model size.}
The raw-RMSE sum gives each task one averaged contribution, but the target scale determines that contribution's magnitude. Large Fibonacci errors dominate many totals, especially at state widths that cannot exactly represent the recurrence. A rank under this measure differs from a rank under median scaled error. The architecture experiment measures the outputs and resource costs of 30 proposals. Its declared operations include a direct product and a hyperbolic tangent activation, so exactness on square and hyperbolic tangent can follow from one primitive. Zero stored coefficients attain the nonnegative coefficient-count lower bound. A smallest graph, smallest total storage or least arithmetic for every task would require additional proofs. Serialized graph storage is reported even when there are no stored coefficients.

\section{Conclusion}
Holding the task and graph fixed makes it possible to measure three different properties of an optimizer's final weights: how well they fit the displayed examples, how close they come to the right answer on new inputs, and whether they reproduce every reference exactly. Across the fixed-network benchmark, iterative WSP has the smallest summed raw extrapolation RMSE among the 14 evaluated pipelines. Its lower total coexists with regressions on individual tasks and an unchanged count of 30 fully exact settings. Allowing Astra to choose the graph yields 21 entirely exact extrapolation sets in 30 further requests, with costs determined by the permitted neural operations. Reporting the returned weights, input/output error and execution cost together exposes these tradeoffs.



